\section{Discussion}
\label{sec:discussion}

\subsection{What is new in the threshold law}

The identity \cref{eq:projected-measure} is stronger than agreement of a
few delay moments: the entire atomic measure seen by the critical mode is
fixed.  A scalar reduction onto that mode is therefore identical along the
\(\eta\)-family.  Nevertheless, \cref{eq:main-law} shows a nonzero
connection displacement.  The missing information is the transverse
measure \cref{eq:transverse-measure}.  It first generates the stable graph
response \cref{eq:H1-eta}; the FitzHugh--Nagumo nonlinearity then returns
that response to the critical field through
\cref{eq:transverse-return-coefficient}.  The Gaussian gap pairing converts
this vector-field return into a scalar root shift.

This separates the result from two familiar observations.  First, the
implicit-function step is not the contribution by itself; its derivative
is meaningful only after a complete-history gap, a one-sided selection, and
the coefficient with a controlled remainder have been constructed.
Second, trajectory agreement in a projected observable would not suffice.
The equivalence \cref{eq:gap-history-equivalence} uses local injectivity on
the common phase section and the injective history embedding, so the root
represents equality in the RFDE history space.

\subsection{Relation to the 2026 delayed van der Pol expansion}

The nearest comparison is \citet{zhang2026high}.  In their scalar delayed
van der Pol equation, weak feedback and long physical delay are scaled as
\(J=O(\eps)\) and \(\tau=O(\eps^{-1/2})\).  A nonlocal center-manifold
reduction and nonlinear time transformation then produce high-order
approximations of the critical parameter and canard curve.  Their scalar
coefficient is an essential calibration for any network model that reduces
exactly to one synchronous delayed oscillator.

The present theorem neither repeats nor imports that scalar coefficient.
It holds the complete critical projected delay operator fixed and varies a
transverse organization that a scalar equation cannot represent.  The new
proof obligation is accordingly at history level: a singular transverse
block is removed uniformly, its mixed parameter jets are extended to a
logarithmic tube, and two one-sided histories are matched.  If the present
model were collapsed to a single oscillator with a uniform self-delay, the
transverse mechanism would disappear and the genuinely new coefficient
\cref{eq:transverse-return-coefficient} would vanish.

\subsection{The scope of the canonical preparation}

The exact root is indexed by \(\cP\).  This is not merely notation.  A
retarded equation has a forward semiflow on history space, and a repelling
one-sided trace cannot be obtained by reversing an arbitrary RFDE initial
history.  Our construction instead chooses both traces relative to explicit
tail levels and realizes their retained pieces in a complete planar flow
inside the invariant history graph.

There is, however, a quantitative transfer principle.  Suppose a separately
constructed physical attracting/repelling history family enters the same
uncut graph and, after the common phase is imposed, its boundary mismatch
and the \(\nu,\eta\)-derivatives used here are bounded by
\[
 C\langle S_\del\rangle^m\e^{cS_\del}.
\]
Then the one-sided Green formulas make the corresponding gap and derivative
errors
\[
 E_\del=O\left(\e^{-S_\del^2/2+cS_\del}
                    \langle S_\del\rangle^{m'}\right)
 =\del^{p-o(1)}.
\]
Since \(\partial_\nu\mathfrak d_\cP\ge c\del\), the root error in the
unfolding coordinate is at most \(CE_\del/\del\), and the physical
parameter error is at most \(C\del E_\del\).  If the same estimate holds
for the \(\eta\)-derivative of the gap mismatch, the error in the
\(\eta\)-displacement law is \(C\del E_\del|\eta|\), below the remainder
in \cref{eq:main-law} for sufficiently large fixed \(p\).

This statement remains conditional for an arbitrary outer Fenichel
selection.  Fixed-parameter closeness is insufficient: without a
parameter-coherent selection rule, an exponentially small boundary
perturbation can have rapidly oscillating parameter derivatives.  A future
outer Lyapunov--Perron construction would have to establish these
full-history derivative bounds rather than infer them from current-state
closeness.

\subsection{Broader implications and next problems}

The mechanism is not tied to the displayed matrices.  In a general
multi-module system, the critical delay measure
\(\ell^T\mathbb B(\dd\theta)r\) records only the direct critical channel.
Unless every delay layer preserves the critical subspace and all nonlinear
blocks respect the same decomposition, a transverse resolvent can feed back
into the critical solvability condition.  The two-layer family here is a
minimal positive example in which this information loss is exact and its
first return can be calculated.

Two extensions are mathematically distinct.  The first is to replace the
local canonical selection by a global physical canard or pulse-onset
selection.  That requires the outer full-history estimate just described.
The second is to pass to a finite heterogeneous network.  There one must
control all transverse delayed modes uniformly and cannot replace the
relevant history operator by the spectral gap of a static adjacency matrix.
Only after that transfer problem is solved would it be justified to combine
the canard threshold with frequency and amplitude as three independently
assigned control coordinates.

Thus the present result supplies one deliberately narrow but reusable
bridge: identical critical projected delay data need not imply identical
nonhyperbolic thresholds, and the discrepancy can be obtained as a rigorous
complete-history root law.  It leaves the biological interpretation of a
global pulse threshold and the design of general network controllers as
separate theorem problems rather than consequences of the local formula.
